% sci_sim_layers.tex — Three simulation layers for the SCI tutorial (S5)
% シミュレーションの3層構造図。SCIチュートリアル S5 用。
%
% Source: docs/architecture/simulation-policy.md §2 (three-layer table) and
% §4 (model-match gate).
% 出典: docs/architecture/simulation-policy.md §2（3層構造の表）と
% §4（モデル一致ゲート）。
%
% Layout: three side-by-side boxes (Layer 1 -> 2 -> 3, left to right) with a
% horizontal arrow chain showing increasing physical/code fidelity — a wide,
% short shape that fits a 16:9 slide without being scaled down to
% illegibility (an earlier stacked-box version was too tall; see git log).
% A gate box spans the full width below, with a feedback arrow on the left
% margin looping the SILS-derived identification back up to Layer 1.
% レイアウト: 3つの箱を左から層1→層2→層3の順に横に並べ、水平の矢印連鎖で
% 物理・コードの忠実度が増していくことを示す。横長で低い形状にすることで、
% 16:9スライドに収めても文字が潰れない（縦積み版は高さが出過ぎたため変更。
% git log 参照）。下に全幅のゲート箱を置き、左マージンのフィードバック矢印で
% SILS側の同定値を層1へ戻すループを示す。
\documentclass[border=10pt]{standalone}
\usepackage{tikz}
\usetikzlibrary{arrows.meta, positioning, calc}
\usepackage{luatexja}
\usepackage[match]{luatexja-fontspec}
% Cross-platform font: Hiragino (macOS) → Yu Gothic (Windows) → Noto (Linux)
% クロスプラットフォーム: ヒラギノ (macOS) → 游ゴシック (Windows) → Noto (Linux)
\usepackage{ifthen}
\newboolean{jfontset}\setboolean{jfontset}{false}
\IfFontExistsTF{Hiragino Kaku Gothic ProN}{%
  \setmainjfont{Hiragino Kaku Gothic ProN}%
  \setsansjfont{Hiragino Kaku Gothic ProN}%
  \setboolean{jfontset}{true}}{}
\ifthenelse{\boolean{jfontset}}{}{%
  \IfFontExistsTF{Yu Gothic}{%
    \setmainjfont{Yu Gothic}%
    \setsansjfont{Yu Gothic}%
    \setboolean{jfontset}{true}}{}}
\ifthenelse{\boolean{jfontset}}{}{%
  \IfFontExistsTF{Noto Sans CJK JP}{%
    \setmainjfont{Noto Sans CJK JP}%
    \setsansjfont{Noto Sans CJK JP}}{}}

\begin{document}

\definecolor{s1color}{RGB}{0, 100, 180}
\definecolor{s2color}{RGB}{40, 140, 60}
\definecolor{s3color}{RGB}{150, 90, 20}
\definecolor{gatecolor}{RGB}{170, 30, 30}
\definecolor{gray1}{RGB}{120, 120, 120}

% Plain-LaTeX helper macros for the colored text lines inside each box
% 箱の中の色付きテキスト用マクロ
\newcommand{\ttier}[2]{{\bfseries\normalsize\color{#1}#2}}
\newcommand{\tsub}[1]{{\scriptsize\color{black!55}#1}}

\begin{tikzpicture}[
    tier/.style={
      rectangle, rounded corners=4pt, draw=#1, line width=1.5pt,
      fill=#1!8, minimum width=60mm, minimum height=30mm,
      align=left, inner xsep=4mm, inner ysep=3mm
    },
    gatebox/.style={
      rectangle, rounded corners=4pt, draw=gatecolor, line width=1.3pt,
      fill=gatecolor!7, minimum width=178mm, minimum height=13mm,
      align=left, inner xsep=6mm, inner ysep=2mm
    },
    arr/.style={-{Stealth[length=3mm]}, line width=1.5pt, gray1},
    garr/.style={-{Stealth[length=3mm]}, line width=1.4pt, gatecolor},
  ]

  % Three layers, left (Layer 1, design-time) to right (Layer 3, SILS)
  % 3層。左（層1・設計用）から右（層3・SILS）まで
  \node[tier=s1color] (S1) at (0, 0) {%
    \begin{minipage}{54mm}\raggedright
      \ttier{s1color}{層1: 設計用線形モデル}\\[3pt]
      \tsub{プラント: $G(s){=}b\,e^{-Ls}/(s(Ts{+}1))$}\\[2pt]
      \tsub{制御コード: 伝達関数として扱う}\\[2pt]
      \tsub{\texttt{sf sysid rate-fit} で同定，設計の正}
    \end{minipage}%
  };
  \node[tier=s2color, right=10mm of S1] (S2) {%
    \begin{minipage}{54mm}\raggedright
      \ttier{s2color}{層2: 実ログ駆動再生}\\[3pt]
      \tsub{プラント: 低次線形 $G(s)$ + 飽和等}\\[2pt]
      \tsub{制御コード: Python 移植，実ログ外乱で駆動}\\[2pt]
      \tsub{\texttt{analysis/scripts/}，A/B 判定の正}
    \end{minipage}%
  };
  \node[tier=s3color, right=10mm of S2] (S3) {%
    \begin{minipage}{54mm}\raggedright
      \ttier{s3color}{層3: SILS}\\[3pt]
      \tsub{プラント: 非線形6自由度（MuJoCo）}\\[2pt]
      \tsub{制御コード: 実ファーム C++ そのもの}\\[2pt]
      \tsub{Code/Param Identity，CI 回帰 33/42 本}
    \end{minipage}%
  };

  % Rightward arrow chain, offset above the boxes (clear of both box border
  % and text — routed entirely in open space, T4 compliant)
  % 右向き矢印の連鎖。箱の上外側の空間に配線
  \draw[arr] ($(S1.north -| S1.east)+(0,6mm)$) -- ($(S2.north -| S2.west)+(0,6mm)$);
  \draw[arr] ($(S2.north -| S2.east)+(0,6mm)$) -- ($(S3.north -| S3.west)+(0,6mm)$);

  % Single label above the whole arrow chain (horizontally centered over S2)
  % 矢印の連鎖全体に対する説明ラベル（層2の真上、中央に1つだけ配置）
  \node[font=\small, text=gray1, align=center, anchor=south]
    at ($(S2.north)+(0,10mm)$) {物理・コードの忠実度が上がる};

  % Model-match gate: a full-width feedback box below all three layers,
  % connected by a return path on the LEFT margin back up to Layer 1
  % (right-angle routing, >=5mm stem before each arrowhead).
  % モデル一致ゲート: 3層の下に全幅のフィードバック箱。左マージンの
  % 折り返し経路で層1へ戻す（直角ルーティング、矢頭前に5mm以上の幹）。
  \node[gatebox, below=20mm of $(S1.south)!0.5!(S3.south)$] (GATE) {%
    \begin{minipage}{170mm}\raggedright
      \ttier{gatecolor}{モデル一致の合否判定}\quad
      \tsub{層3のログに同一の同定パイプラインを適用し $(b,L,T)$ を層1の実機基準値と比較（\texttt{sf sils sysid-gate}）}
    \end{minipage}%
  };

  % NOTE: the vertical bus MUST use a single x-coordinate anchored to S1
  % (not to GATE.west), because GATE is centered on the S1-S3 span and its
  % west edge is well to the right of S1's west edge — using GATE.west's own
  % x for the bus would route the line back through S1's interior (T4
  % violation). Anchoring both turns to feedX keeps the whole bus strictly
  % in the left margin, clear of every box.
  % 注意: 垂直バスは必ず S1 基準の単一 x 座標を使うこと（GATE.west 基準にしない）。
  % GATE は S1〜S3 の全幅中央に配置されるため GATE.west は S1.west よりかなり
  % 右にある。GATE.west 基準で配線すると S1 内部を貫通する（T4違反）。
  % feedX に両方の折れ点を揃えることで、バス全体を左マージンに収める。
  \coordinate (feedX) at ($(S1.west)+(-8mm,0)$);
  % Fix (T4): "-|" takes the x-coordinate from its RIGHT operand and the
  % y-coordinate from its LEFT operand, so the corner must be written as
  % (GATE.west -| feedX) = (feedX's x, GATE.west's y). The previous
  % (feedX -| GATE.west) = (GATE.west's x, feedX's y) instead put the corner
  % at GATE.west's x (~-19mm), which lies inside S1's x-range ([-30,30]mm)
  % and drove the bus straight through the S1 box.
  % 修正（T4）: 「-|」は x 座標を右側の被演算子から、y 座標を左側の
  % 被演算子から取る。よってコーナーは (GATE.west -| feedX) =
  % (feedXのx, GATE.westのy) と書く必要がある。以前の
  % (feedX -| GATE.west) は (GATE.westのx, feedXのy) となり、コーナーが
  % GATE.westのx（約-19mm）になっていた。これはS1のx範囲（[-30,30]mm）の
  % 内側であり、バスがS1の箱を貫通していた。
  \draw[gatecolor, line width=1.4pt]
    (GATE.west) -- (GATE.west -| feedX) -- (feedX);
  \draw[garr] (feedX) -- (S1.west);
  % Label sits clear to the LEFT of the vertical bus/arrow corner (feedX).
  % NOTE: the midpoint must be captured as its own coordinate first — TikZ
  % calc parses "(A)!0.5!(B) + (dx,dy)" by attaching "+(dx,dy)" to B, not to
  % the interpolated result, which silently halved the intended x-offset in
  % an earlier revision (label barely moved, still grazed by the arrow).
  % Widening the offset to -8mm and the S1-GATE gap (12mm -> 20mm, below)
  % together keep the rotated text's right edge and vertical span clear of
  % the arrow's corner and arrowhead. Previously the arrowhead crossed the
  % "と" character (T2).
  % ラベルは垂直バス/コーナー(feedX)から左に十分離して配置する。
  % 注意: 中間点は必ず独立した座標として確定させること — TikZ calc は
  % 「(A)!0.5!(B) + (dx,dy)」を「B に +(dx,dy)」と解釈するため、以前の版では
  % 意図したxオフセットが半分になり、ラベルがほぼ動かず矢印に接触したまま
  % だった。オフセットを-8mmへ広げ、層1〜3とGATEの間隔（below）も12mm→20mm
  % に広げたことで、回転テキストの右端と縦方向の占有幅の両方が、矢印の
  % コーナー・矢頭から確実に離れる。以前は矢頭が「と」の文字に重なって
  % いた（T2違反）。
  % Use the same corrected corner point as the bus above so the label sits
  % on the actual vertical segment (in the left margin), not on the old
  % (now-removed) mis-routed segment.
  % 上のバスと同じ修正後のコーナー点を使う。ラベルを実際の垂直区間
  % （左マージン側）の上に置くため（以前の誤配線区間の上ではなく）。
  \coordinate (feedMid) at ($(feedX)!0.5!(GATE.west -| feedX)$);
  \node[font=\scriptsize, text=gatecolor, align=center, rotate=90]
    at ($(feedMid) + (-8mm,0)$) {実機同定値と突き合わせ};

\end{tikzpicture}
\end{document}
